
\subsection{Projective gauge incidence of the common constructor and descent to the quotient}

The gauge image is not the link marginal. Its complete object at each level is
\begin{equation}\label{eq:pdf2-g-gau-explicito}
 \mathfrak G_{{\rm gau},m}(d_{\rm gau})=
 \bigl(
 \begin{gathered}
  \mathfrak G_{\rm gau}^{\rm int};\\
  \Lambda_m,\widehat X_m,\widehat\mu_m^{\rm ov},
  \{\widehat\mu_{m,g}^{\rm YM},\nu_{m,g},K_{m,g}^{\rm phys}\}_{g>0},\\
  \mathcal G_m^{\rm full},\mathbb P_m^{\rm phys},
  \widehat J_m,\widehat\Gamma_m,
  \{E_{m,c}\}_c,\widehat H_m^{\rm ov},D_m^{\rm YM},\\
  \{\Theta_{\nu,m}\}_{\nu=1}^4,
  \{\mathfrak A_{+,\nu,m},\operatorname{ev}_{t}^{(\nu)},
      \mathcal V_{m,\nu}^{\rm OS}\}_{\nu=1}^4,\\
  \widehat\Omega_m^{(8)},P_{\Omega_m},W_c,
  \operatorname{Carr}_m,\operatorname{Mem}_m,
  \operatorname{Surv}_m,\tau_m,\mathcal P_m^{F^2}
 \end{gathered}
 \bigr).
\end{equation}
The group $\mathcal G_m^{\rm full}$ acts simultaneously on links, frames, and
incidence coordinates $q_c\in\mathbb F_3^6$. The physical projection is its Haar
average. The reduced generator is
\begin{equation}\label{eq:pdf2-g-gau-generador-estructural}
 D_m^{\rm YM}
 =\left.3\kappa_{\rm av}\sum_c(I-E_{m,c})
  \right|_{\operatorname{Ran}\mathbb P_m^{\rm phys}},
\end{equation}
and $W_c$ belongs to this subalgebra. The four reflections act on the same
measure, the same kernel, and the same terminal.
The preceding exact source construction establishes the incidence gauge action
\eqref{eq:amp-ym-accion-gauge-incidencia}, the physical Haar projection
\eqref{eq:amp-ym-proyeccion-fisica-haar}, and the intertwined compression
\eqref{eq:amp-ym-compresion-entrelazada}. This construction yields the
complete equivariant kernel $\widehat K_{m,t}^{\rm ov}$, its descent to the
physical quotient $K_{m,t}^{\rm phys}$, the four reflection operators, and
the action $U_m(g)$ of $E(4)$ on the smooth cylindrical core. This block
assembles, without introducing them as data, the generator, vacuum, curvature
witness, and OS forms:
\begin{equation}\label{eq:pdf2-g-gau-entrelazamiento-dato}
 D_{m+1}^{\rm YM}\widehat J_m=\widehat J_mD_m^{\rm YM},
 \quad
 \mathbb P_{m+1}^{\rm phys}\widehat J_m
  =\widehat J_m\mathbb P_m^{\rm phys},
 \quad
 \widehat J_mW_{m,c}=W_{m+1,c}.
\end{equation}
The first and third identities express, in the common
notation, \eqref{eq:amp-ym-compresion-entrelazada} and
\eqref{eq:amp-ym-testigo-fisico-gap}; the semigroup, simple vacuum, and
exact gap have already been obtained in
\eqref{eq:amp-ym-owner-semigrupo-vacio}--
\eqref{eq:amp-ym-owner-brecha-exacta}. Here
$\nu_{m,g}:=(\pi_m^{\rm phys})_*\widehat\mu_{m,g}^{\rm YM}$ and
$K_{m,g}^{\rm phys}(t)=e^{-tD_{m,g}^{\rm YM}}$ on the physical quotient.
Gauge completeness also includes disintegration with respect to each collection of
Euclidean hyperplanes. If $A_j\in\mathfrak A_{+,\nu,m}$ and
$t_1\le\cdots\le t_n$, the marginal of the same physical measure satisfies
\begin{equation}\label{eq:pdf2-g-gau-transfer-euclidean-dato}
\begin{aligned}
 &\int\prod_{j=1}^n
   (\tau_{t_je_\nu}A_j)(x)\,d\widehat\mu_{m,g}^{\rm YM}(x)\\
 &\quad=
 \int\prod_{j=1}^nA_j(y_j)\,
  \nu_{m,g}(dy_1)
  \prod_{j=1}^{n-1}K_{m,g}^{\rm phys}
       (t_{j+1}-t_j;y_j,dy_{j+1})\\
 &\quad=
 \left\langle\mathbf1,
  M_{A_1}e^{-(t_2-t_1)D_{m,g}^{\rm YM}}M_{A_2}\cdots
  e^{-(t_n-t_{n-1})D_{m,g}^{\rm YM}}M_{A_n}\mathbf1
 \right\rangle.
\end{aligned}
\end{equation}
The preceding reconstruction proves the identity in all four directions
through \eqref{eq:amp-ym-cuatro-os} and constructs a unique intertwined
Hamiltonian in \eqref{eq:amp-ym-hamiltoniano-comun}. Consequently, OS
translation is expressed here as
\begin{equation}\label{eq:pdf2-g-gau-os-generator-dato}
 \mathcal V_{m,\nu}^{\rm OS}T_{m,\nu}^{\rm OS}(s)
 (\mathcal V_{m,\nu}^{\rm OS})^{-1}
 =e^{-sD_{m,g}^{\rm YM}}.
\end{equation}
Thus, the expectation generator and the translation Hamiltonian are not
identified merely by name: they are two already constructed realizations of the same
disintegration, whose limiting form and curvature output are given in
\eqref{eq:amp-ym-forma-limite} and \eqref{eq:amp-ym-f2}.
